\usepackage{polyglossia} % языковой пакет

\usepackage{pdfpages} % пакет для импорта pdf-файлов

\usepackage{tocvsec2} %%%%%%%%%%%%%%%%%%%%%%%%%%%%%%%%%

% Подписи рисунков и таблиц по требованию кафедры АСУ:
% «Рисунок 1 — Описание», разделитель — длинное тире.
\usepackage[labelsep=endash]{caption}
% Polyglossia сбрасывает \figurename на «Рис.» при активации русского
% языка, поэтому \renewcommand{\figurename}{...} в преамбуле и даже
% \AtBeginDocument срабатывают ненадёжно. Самый прямой способ —
% задать имя метки в caption-пакете через name=..., минуя \figurename.
\captionsetup[table]    {name=Таблица, labelsep=endash, justification=raggedleft, singlelinecheck=off, position=above}
\captionsetup[figure]   {name=Рисунок, labelsep=endash, justification=centering,  singlelinecheck=on,  position=below}
\captionsetup[lstlisting]{labelsep=endash, justification=centering, singlelinecheck=on,  position=below}
% Дублируем через \AtBeginDocument для пакетов, читающих \figurename
% напрямую (например, hyperref для подписей якорей).
\AtBeginDocument{%
    \renewcommand{\figurename}{Рисунок}%
    \renewcommand{\tablename}{Таблица}%
}

\usepackage{titletoc}

% \usepackage{misccorr} отключён: пакет добавляет точку после номера
% в подписях рисунков/таблиц (делает «Рисунок 1.1.», «Таблица 1.1.»),
% что мешает применению нашего labelsep=endash. Точки в оглавлении
% и так формируются дотлидером \titlecontents, без misccorr.

\usepackage{graphicx} % пакет для использования графики (чтобы вставлять рисунки, фотографии и пр.)

\usepackage{amsmath} % поддержка математических символов

\usepackage{url} % поддержка url-ссылок

\usepackage{lipsum}

% Legacy команда \bibliographystyle{gost-numeric.bbx} удалена:
% biblatex + biber не используют bibtex'овский bibliographystyle.
% Минимальный набор опций ГОСТ-стиля; deprecated parentracker
% и проблемная пара defernumbers+sorting=none убраны — они вызывали
% Missing number в \printbibliography.
\usepackage[
    backend=biber,
    bibencoding=utf8,
    style=gost-numeric,
    language=auto,
    autolang=other,
    sorting=nty
]{biblatex}
% Принудительная сортировка: сначала русскоязычные источники,
% затем иностранные. Реализуется через поле presort в refs.bib:
% русские записи — presort={a}, латинские — presort={z}.
% При sorting=nty biber сортирует presort → name → year → title.
\addbibresource{refs.bib}

% Переопределение биб-локализации под ГОСТ Р 7.0.100-2018:
% - urlseen ('дата обращения' вместо biblatex-gost-дефолта 'дата обр.')
% - bibrangedash (короткое тире en dash в диапазонах страниц '23–31'
%   вместо длинного тире, которое biblatex-gost подставляет в русском
%   режиме).
\DefineBibliographyStrings{russian}{%
    urlseen = {дата обращения},
}
\AtBeginDocument{\def\bibrangedash{\textendash}}

\usepackage{fancyhdr} % заголовок

\usepackage{enumitem} % списки

\usepackage{multirow} % таблицы с объединенными строками

\usepackage{hyperref} % пакет для интеграции гиперссылок

\usepackage{indentfirst} % пакет для отступа абзаца


\usepackage{chngcntr} % пакет подписей и нумерации рисунков

\usepackage{setspace}

\usepackage{csquotes}

\usepackage{tabularx}

% Альбомная ориентация для широких диаграмм (use case, flowchart).
\usepackage{pdflscape}
\usepackage{rotating}

% Барьер для плавающих объектов (рисунков и таблиц): \FloatBarrier
% не позволяет рисунку «уплывать» в следующий подраздел.
\usepackage{placeins}
% Пакет float даёт спецификатор [H] для строгой фиксации рисунка
% в указанном месте — используется в макросе \imageH (без float).
\usepackage{float}

% Эмуляция пакета needspace (пакет в TL2026 basic недоступен
% без прав администратора). \needspace{N\baselineskip} ломает
% страницу, если на текущей странице осталось места меньше N.
% Используется для предотвращения висячих заголовков перед
% подразделами и листингами в приложениях.
\makeatletter
\providecommand\needspace[1]{%
    \par\penalty-100%
    \begingroup
        \dimen@\pagegoal
        \advance\dimen@-\pagetotal
        \ifdim\dimen@<#1\relax
            \pagebreak
        \fi
    \endgroup}
\makeatother

% Параметры выключки абзаца: уменьшаем «реки» пробелов в узких
% колонках русского текста за счёт более гибкого emergencystretch
% и более терпимого hyphenpenalty при выключке по ширине.
\tolerance=1000
\emergencystretch=3em
\hyphenpenalty=50

%%%%%%%%%%%%%%%%%%%%%%%%%%%%%%%%%%%%%%%%%%%%%%%%%%%%%%%%%%%%%%